\documentclass[12pt]{article}

% Packages
\usepackage[margin=1in]{geometry}
\usepackage{setspace}
\usepackage{amsmath}
\usepackage{amssymb}
\usepackage{graphicx}
\usepackage{xcolor}
\usepackage{booktabs}
\usepackage{array}
\usepackage{tabularx}
\usepackage{hyperref}
\usepackage{apacite}
\usepackage{natbib}

% Title information
\title{Optimal Flood Policy}
\author{Taco Prins, Lennard Schlattmann, and Daniel Jonas Schmidt}
\date{\today}

\begin{document}

\maketitle

\section{Motivation}

%Flooding is the most common natural disaster in the United States, and the costliest in annual damages [reference]. Damages from flooding are projected to increase due to climate change and population movements towards exposed areas ([reference];\cite{wing2022inequitable}). 

Flooding is the most common and costliest natural hazard faced by U.S. homeowners, with losses to homeowners with federally backed mortgages alone averaging about $\$$9.4 billion per year under current climate conditions (\cite{cbo2023flooddamage}). A significant increase in flood risk is projected under a wide range of climate scenarios, and adaptation by homeowners to this risk has been studied in economics as a key variable mediating the resulting increase in damages (\cite{desmet2018evaluating}). 

Historically, almost all protection from flood risk enjoyed by homeowners has been provided by federal government programmes. These seek to reduce the total amount of risk, while also transferring costs from homeowners to taxpayers. Direct government assistance to homeowners in floodprone areas takes two main forms: a public flood risk insurance offer known as the National Flood Insurance Programme (NFIP) and ex-post disaster aid provided to affected households by the Federal Emergency Management Agency (FEMA). Through a mixture of incentives and regulation, the federal government furthermore encourages the building of more resilient housing in the most floodprone areas.

A key motivation of such policy interventions is that outcomes in an unfettered private market would leave homeowners exposed to excessive risks. The Act of Congress establishing the NFIP in 1968 noted that `many factors have made it uneconomic for the private insurance industry alone to make flood insurance available to those in need of such protection on reasonable terms and conditions' (\cite{nfia1968}). Another problem, however, is that U.S. homeowners do not always fully take flood risk into account. This fact is very widely documented and is evident both in survey evidence (e.g. \cite{bakkensen2022going}; \cite{mulder2024mismeasuring}) and in market outcomes regarding house prices, flood insurance uptake, and flood adaptation rates. 

For example, the price discount of homes exposed to flood risk is found to be less than the discounted value of expected damages in many empirical studies\footnote{\cite{fairweather2024expecting} cite a list of 20 such studies.}; homebuyers with climate sceptic beliefs pay more for houses exposed to future flood risk through sea level rise than other buyers (\cite{bernstein2019disaster}; \cite{baldauf2020does}; \cite{bakkensen2025leveraging}); and the discount for flood-exposed housing waxes and wanes over time as current events influence the salience of the risk in   homebuyers' minds (\cite{keys2020neglected}; \cite{giglio2021rfs}). \cite{wagner2022adaptation} and \cite{amornsiripanitch2025measuring} note that the uptake of flood risk insurance in high-risk Special Flood Hazard Areas (SFHAs) is relatively low (not more than 60$\%$ in their samples), even though insurance in these areas is mandated for many homeowners and for many properties premiums have historically been less than expected annual losses.\footnote{Furthermore, \cite{gallagher2014learning} uncovers a dynamic of Bayesian updating with forgetting that explains time-varying flood risk insurance demand at the county level as a function of nearby floods, similar to the observed salience effect of floods on house prices.} Finally, the recent empirical study of \cite{ostriker2022effects} finds that houses adapted to flooding in Florida do not command a detectable price premium, and publicly available data released by FEMA\footnote{LINK} reveal that effective adaptation measures are rare outside of areas where they are legally required.




In this paper, we formalise the trade-off faced by a policymaker dividing a fixed budget between a premium subsidy for flood risk insurance and ex-post disaster aid. We derive simple formulas for the marginal welfare gain from an extra dollar spent on either policy instrument. [1-2 further sentences describing more work.]

Some terms in our welfare formulas evaluate the transfers to different households inherent in these policies: holding insurance demand constant,  premium subsidies redistribute money to insured households, and disaster aid to uninsured households affected by floods. Households receiving disaster aid on average have higher marginal utility of consumption because floods destroy wealth and because uninsured households are poorer.

The remaining terms arise from the behavioural response to the policy instruments. First,  more generous disaster aid reduces households' use of subsidised insurance and vice versa. There is a fiscal spillover between disaster aid and premium subsidies as part of the cost of one policy is recouped through reduced spending on the other. Second, when households that underestimate flood risk are indifferent between buying insurance or not, a policy change that induces them to buy makes them better off.\footnote{The household's utility then increases by the difference between its subjective welfare estimate, conditional on being uninsured, and an objective estimate calculated using the true probability of flooding.}  This welfare gain appears in our formulas as a so-called behevioural internality. To quantify it, we posit a distribution of household beliefs over flood risk. Unlike the fiscal spillover, which depends only on the total change in insurance demand, the size of the internality follows from the location in the belief distribution where the change in insurance demand occurs.

Our key finding is that, over the range of plausible calibrations, optimal spending on disaster aid as a proportion of the total flood policy budget is higher than its current level. Although cumulative funding allocated to the NFIP programme has far exceeded the amount of disaster aid granted by FEMA, a substantial fraction of homeowners do not perceive flooding to be a sufficiently significant risk to purchase insurance even at a subsidised rate. This leaves disaster aid as the major safety net for this group. Moreover, the moral hazard effect of increased disaster aid is weaker when flood risk beliefs are subjective. The intuition is that households far away from the point where willingness to pay for insurance equals the actual premium are not responsive to marginal changes to the cost or benefit of insurance. Errors in flood risk perception thus decrease the sensitivity of insurance demand to both premium subsidies and disaster aid.    

In our baseline model, households vary only in their flood risk beliefs and buy insurance if their subjective flood arrival rate exceeds a threshold value. We assign values to components of the formula regarding households preferences and beliefs, the process for flood risk, and the behavioural response to policy changes to match estimates in the empirical literature. The optimal proportion of uninsured flood damages compensated by disaster aid then exceeds its current level of around 9$\%$\footnote{This number is taken from \cite{gruber2026optimal}, using their IV estimate of $\$215.70$ in disaster aid spending per flood-affected house, dividing by the application rate for disaster aid of $8\%$, and dividing the total again by the estimated average damage per house of $\$29,267$. Their definition of disaster aid encompasses not only the IHP, SBA, and HMGP programmes but also defaults on GSE-backed mortgages.} to reach x$\%$. The optimal premium subsidy amounts to x$\%$ of the actuarially fair premium. This result runs counter to recent reforms to the NFIP programme that mostly eliminate premium subsidies, but is smaller than the average subsidy that pertained for most of the programme's history.   

In subsequent extensions, we investigate the robustness of our central result. First, we examine the redistributive effect of expanded disaster aid across the income distribution. Second, we consider its effect on the demand for adapted housing. Third, we introduce variation in risk aversion as a further source of household heterogeneity. The qualitative result that optimal disaster aid as a fraction of flood damages exceeds its current level holds up in each extension. We also note that accounting for differences in administrative costs strengthens the case for disaster aid as a larger part of the policy mix.     

Our study is most closely related to the recent contribution by \cite{gruber2026optimal}, which analyses optimal premium subsidies holding fixed the amount of disaster aid. There, the key formula which is solved by the optimal premium subsidy is composed of empirically observable elasticities. The central estimate of the optimal flood risk insurance premium subsidy turns out to be 52$\%$, not far from the average subsidy of around 50$\%$ which according to the authors applied before the recent RR2.0 reform of the NFIP programme. Part of what renders the subsidy worthwhile is its role in dampening the reclassification risk that comes from risk-based insurance pricing. Even conditional on the amount of global atmospheric warming, projected flood risk varies greatly between climate models, and, with actuarially fair pricing, so do insurance premiums. Premium subsidies absorb the demeaned property-level component of this risk which homeowners otherwise cannot insure against.

The quantitative results of \cite{gruber2026optimal} are sensitive to two issues. First, in an extension where they allow for underestimation of flood risk, the optimal subsidy rises to 89$\%$, implying that the government almost fully absorbs flood risk faced by homeowners. Second, in their framework, premium subsidies and disaster aid are strongly complementary as eliminating FEMA disaster aid lowers the optimal subsidy from 52$\%$ to 15$\%$. As they write, this `joint change is welfare-improving only if households are also making rational risk and FEMA-generosity decisions' (ibid p.35), which they admit is unlikely to be the case. Importantly, we solve for  optimal public spending on premium subsidies and disaster aid without assuming either.

This paper proceeds in four parts. Section 2 briefly describes the institutional setting determining current flood insurance provision and disaster aid spending in the United States. Section 3 introduces our model of household decision-making and welfare and states the key formulas solved by the optimal levels of the premium subsidy and disaster aid. Section 4 discusses various threats to the robustness of our result and develops model extensions to address some of them. Section 5 concludes and offers suggestions for further research.    

% As annual damages are projected to increase due to climate change, as well as population movements to vulnerable areas [references], the financial sustainability of government programmes has received the attention of an expanding economic literature.   


\section{Institutional setting}

\subsection{Insurance}
The NFIP programme ran a deficit almost immediately upon its establishment in 1968 as on average premiums were less than expected flood damages. The impact of Hurricane Katrina in 2005 effectively rendered the programme insolvent (\cite{michel2010catastrophe}). The programme has also been marked by cross-subsidies from low-risk properties to high-risk properties, as premiums were set on the basis of coarse flood risk maps that did not reflect flood risk at the property level (\cite{kousky2018financing}). Two recent reforms, the Biggert-Waters Act of 2012 and the Risk Rating 2.0 (RR2.0) reform of 2021-2022, have sought to eliminate insurance premium subsidies. Especially as a result of RR2.0, premiums have been rising steeply in parts of the United States. Regulation  caps yearly NFIP premium increases at 18$\%$. While around 34$\%$ of policies were at the actuarially fair premium when RR2.0 was introduced, around 30$\%$ were set on a path towards the doubling of premiums or more (\cite{gao2023flood}, p. 25).  

Homeowners with a GSE-backed mortgage inside Special Flood Hazard Areas (SFHAs) are obliged to hold a flood insurance policy. The mandate is imperfectly enforced, however. \cite{gruber2026optimal} estimate that around 18$\%$ of all homeowners are insured when weighted by flood damages, which includes both homeowners inside and outside SFHAs. Insurance take-up is higher in SFHAs at around 50-60$\%$, but the majority of average annual losses are suffered by homeowners outside these areas \cite{amornsiripanitch2025measuring}.

Although the NFIP provides homeowners with the opportunity to insure at an actuarially fair premium, it has coverage limits of $\$$250,000 for property structures and $\$$100,000 for personal contents. These limits are occasionally binding, as shown by \cite{sastry2026bears}.\footnote{As shown in Figure 1.2 of the paper's appendix, the 90th percentile of claims between 2008-2018 in the states of Florida, Texas, Alabama, Mississippi, and Louisiana  levels off at the $\$$250,000 limit for homes with an estimated replacement cost exceeding $\$$400,000 (panel (a)), which form a non-trivial part of the distribution (panel (b)).} Private insurance policies, which are generally taken out in combination with an NFIP policy by owners of higher-value properties as supplementary insurance, accounted for roughly 3.5 to 4.5$\%$ of all primary residential flood policies in 2018 (\cite{kousky2018emerging}).

\subsection{Disaster aid}
When the president issues a Presidential Disaster Declaration (PDD)\footnote{Formally, the president may declare either an emergency or a major disaster under the Stafford Act, with the former releasing less funds. The SBA Administrator may also separately issue a declaration that makes SBA disaster loans available.} after a flood, individuals may apply for two main forms of so-called Individual Assistance.\footnote{Following a disaster declaration, FEMA also oversees direct transfers to states and local governments within designated counties. The Public Assistance (PA) programme covers emergency measures and repairs, while the Hazard Mitigation Grant Programme (HMGP) finances projects that reduce the costs from future disasters.} First, FEMA's Individuals and Household Program (IHP) consists of grants which cover repair expenses and other immediate housing costs (Housing Assistance) and non-housing costs (Other Needs Assistance). Intended to meet `basic needs and support recovery efforts' (\cite{crs2024disaster}), the maximum size of each component was $\$$44,800 as of 2026, while the average IHP grant between 2002-2024 was $\$$3,446 \cite{crs2025ihpfaq}. Second, the Small Business Administration (SBA), not directly administered by FEMA, provides loans at subsidised rates\footnote{The annual interest rate of SBA loans may not exceed 4$\%$ `if the applicant is found 
by SBA to be unable to obtain credit elsewhere' (\cite{crs2024disaster}, p.2).} with maturities of up to 30 years. These loans are capped at $\$$500,000 for home repairs and $\$$100,000 for personal property replacement. 

Homeowners with GSE-backed mortgages form a special category of aid recipients. A disaster declaration that releases Individual Assistance funding also makes homeowners eligible for forbearance for up to 12 months on their mortgage repayments for GSE-backed loans.\footnote{Information accessed via this \hyperlink{https://www.fhfa.gov/homeowners-and-homebuyers/mortgage-assistance/disaster-assistance}{link}.} Furthermore, it has been argued that mortgage lenders' ability to sell securitised mortgages to GSEs encourages them not to adjust loan conditions to flood-related default risk, giving homeowners access to underpriced implicit insurance against flood damages  (\cite{keenan2020underwaterwriting}; \cite{sastry2023insurers}). \cite{bakkensen2025leveraging} and \cite{liaomulder2026homeequity} find evidence that homeowners exploit this insurance option in a strategic manner.\footnote{\cite{bakkensen2025leveraging} show that buyers identified as pessimistic regarding climate change are more likely to use a mortgage when purchasing properties exposed to sea level rise and have higher loan-to-value ratios,. \cite{liaomulder2026homeequity} show that owners of highly leveraged properties are less likely to purchase flood risk insurance.} However, \cite{gallagher2017household} note that Hurricane Katrina engendered at most a modest rise in defaults, while \cite{sastry2026bears} demonstrates that mortgage lenders restrict credit in flood-prone areas to avoid excess losses, mainly by reducing maximum loan-to-value ratios.\footnote{\cite{nguyen2022climate} provide earlier evidence of reduced LTVs in flood-prone areas, but, using data , note that providers underestimate.}\cite{qin2026flood}). 


\subsection{Adaptation}

Notes on subsidies and regulations encouraging regulation, and on public flood protection programmes (covering 1.5$\%$ of homeowners, \cite{qin2026flood}).

\bibliographystyle{apacite}
\bibliography{bibfile.bib}
\end{document}